\begin{trapbox}
	\begin{description}[style=nextline, leftmargin=2.8em, labelsep=0.5em, itemsep=0.35em]
		\item[\textbf{\textcolor{CTrap}{易错点一（线段混淆）}}] 误认为 $PT^{2}=PA\cdot AB$，即错误地将 $AB$ 作为乘积中的一段。
		\item[\textbf{\textcolor{CTrap}{正确理解（割线两段区分）：}}] 乘积应是 $PA\cdot PB$，即点 $P$ 到两个交点的线段长之积。
		\item[\textbf{\textcolor{CTrap}{错因分析（位置关系不清）：}}] 未能正确识别 $A$ 在 $P$、$B$ 之间，误将 $AB$ 当作 $PB$。
	\end{description}

	\medskip
	\begin{description}[style=nextline, leftmargin=2.8em, labelsep=0.5em, itemsep=0.35em]
		\item[\textbf{\textcolor{CTrap}{易错点二（条件忽视）}}] 在没有明确点 $P$ 在圆外的情况下直接使用切割线定理。
		\item[\textbf{\textcolor{CTrap}{正确理解（外点前提）：}}] 只有 $P$ 在圆外时切线才存在，定理才适用。
		\item[\textbf{\textcolor{CTrap}{错因分析（前提漏检）：}}] 忽略了定理的前提条件，生搬硬套公式。
	\end{description}

	\medskip
	\begin{description}[style=nextline, leftmargin=2.8em, labelsep=0.5em, itemsep=0.35em]
		\item[\textbf{\textcolor{CTrap}{易错点三（顺序颠倒）}}] 当 $A$ 不在 $P$、$B$ 之间时，直接套用 $PT^{2}=PA\cdot PB$ 导致错误。
		\item[\textbf{\textcolor{CTrap}{正确理解（A在P、B之间）：}}] 定理要求 $A$ 是割线上靠近 $P$ 的交点，$B$ 是较远的交点。
		\item[\textbf{\textcolor{CTrap}{错因分析（顺序不注意）：}}] 没有注意点 $A$、$B$ 在割线上的位置顺序。
	\end{description}
\end{trapbox}
